\documentclass[10pt,a4paper]{amsart}
\usepackage[margin=19mm]{geometry}
\usepackage{amsmath,amssymb,mathtools}
\usepackage[scr=rsfs,cal=euler]{mathalfa}
\usepackage{xfrac}
\usepackage[unicode,hidelinks,bookmarks=false]{hyperref}
\newcommand{\ie}{i.e.\ }
\newcommand{\cf}{cf.\ }
\newcommand{\resp}{respectively\ }
\newcommand{\Subsection}{\subsection}
\providecommand{\nobreakdash}{\nobreak\hbox{-}}
\setlength{\emergencystretch}{2em}
\multlinegap0pt
\usepackage{bm}
\usepackage{bbm}
\usepackage{dsfont}
\usepackage{xspace}
\usepackage{cancel}
\usepackage{mathtools}
\usepackage{amsthm}
\makeatletter
\@ifundefined{theorem}{\newtheorem{theorem}{Theorem}}{}
\@ifundefined{lemma}{\newtheorem{lemma}[theorem]{Lemma}}{}
\makeatother
\title[Bound-Preserving WENO Schemes for Temple-class systems]{Bound-Preserving WENO Schemes for Temple-class systems}
\author{Wei Chen, Shumo Cui, Kailiang Wu, Tao Xiong, Baoyue Yu}
\date{Source: arXiv:2510.04123v1 (CC BY 4.0)}
\begin{document}
\maketitle

\noindent\emph{Source and licence.} Bound-Preserving WENO Schemes for Temple-class systems. Version arXiv:2510.04123v1 (CC BY 4.0), distributed under the Creative Commons Attribution 4.0 International licence, as recorded in the arXiv licence metadata for this exact version and shown on the abstract page https://arxiv.org/abs/2510.04123v1. Full rights notice: licenses/arxiv-2510.04123-CC-BY-4.0.txt.\par\smallskip

\noindent\textit{Editorial scope.} Counted unit: the Theorem whose source label is \texttt{infeasible} in the article (source0.tex lines 521--524, source label \texttt{infeasible}), delivered together with the article's own complete original proof of it (source0.tex lines 525--564, one \texttt{proof} environment of 40 lines). No mathematics is written, repaired or re-derived: statement and proof are the article's own text line for line. Required context retained uncounted: newsystem (lines 451--453); 778 (lines 480--483). Every definition, display and label of those lines is retained unchanged. Omitted from this package and not redistributed: 1--450, 454--479, 484--520, 565--1454. Nothing inside a retained excerpt is omitted or altered: every retained body is delivered exactly as the article's lines are written, so no omission receipt of any kind accompanies this package. Citations: the retained excerpts carry no citation command at all, so no bibliography entry of the article is transcribed and no reference is redistributed. Reference adaptation: none was needed and none was made; every reference inside the delivered unit resolves inside it, so no unresolved reference or citation is produced. Printed slips kept literal: no printed word, symbol, hypothesis or number of the counted statement or of its proof was altered. Layout and class: the article's own class is \texttt{elsarticle} with packages including jcomp, framed, multirow, hyperref, algpseudocode, amssymb, amsthm, amsmath, bm, amsfonts, dsfont, latexsym, subcaption, url; the wrapper is the lane's pinned \texttt{amsart} editorial wrapper at 10pt on A4, which restates the theorem-like environments the retained text needs and adds the article's own macro definitions that the retained text needs (none), with extra wrapper packages (bm, bbm, dsfont, xspace, cancel, mathtools). The delivered numbering is the unit's own. Assets: no retained span contains a figure environment, a graphics-inclusion command, a PDF-page inclusion, a table or a TikZ picture; no image file is copied into this package, the package declares no asset dependency and no image file is redistributed with it. Licence: version arXiv:2510.04123v1 of this article is distributed under the Creative Commons Attribution 4.0 International licence, as recorded in the arXiv licence metadata for this exact version and shown on the abstract page https://arxiv.org/abs/2510.04123v1. A full rights notice is delivered at licenses/arxiv-2510.04123-CC-BY-4.0.txt. Characters: every retained excerpt is pure ASCII, so the delivered engine, XeLaTeX with its default Latin Modern text font, reports no missing character.
\begin{equation}\label{newsystem}
    \partial_\tau \boldsymbol{\mathcal{U}} + \partial_\xi \mathbf{G}(\boldsymbol{\mathcal{U}}, c) = 0,
\end{equation}

\begin{lemma}
\label{778}
	The invariant domains $\Omega_1$ and $\Omega_2$ are convex sets, while $\Omega$ is nonconvex.
\end{lemma}

\begin{theorem}
\label{infeasible}
	For the system \eqref{newsystem}, it is impossible to construct a stable BP numerical scheme when $c \neq v$.
\end{theorem}

\begin{proof}
	The numerical scheme for system \eqref{newsystem} can be described by:
	\begin{equation*}
		\boldsymbol{\mathcal{U}}_j^{n + 1} = \boldsymbol{\mathcal{U}}_j^n - \frac{\Delta \tau_n}{\Delta \xi} \left( \widehat{\mathbf{G}}_{j + \frac{1}{2}}^n - \widehat{\mathbf{G}}_{j - \frac{1}{2}}^n \right),
	\end{equation*}
	where $\widehat{\mathbf{G}}_{j \pm \frac{1}{2}}^n$ represents the numerical flux. Note that $\boldsymbol{\mathcal{U}}_j^{n + 1}$ can be regarded as a function of $\Delta \tau_n$, namely, $\boldsymbol{\mathcal{U}}_j^{n + 1} (\Delta \tau_n)$, if the computation of numerical flux has been fixed. Given that $c \neq v$, we can construct a Riemann initial data
	\begin{equation*}
		\boldsymbol{\mathcal{U}}^0 (\xi) = \left\{
		\begin{aligned}
			& \boldsymbol{\mathcal{U}}_L \quad {\rm for} \; \xi \leq \xi_{S + \frac{1}{2}}, \\
			& \boldsymbol{\mathcal{U}}_R \quad {\rm otherwise},
		\end{aligned}
		\right.
	\end{equation*}
	satisfies
	\begin{equation}
	\label{90}
		v \left(\boldsymbol{\mathcal{U}}^0 \right) \equiv v_{\min} \equiv v_{\max}, \quad \frac{\widehat{\mathbf{G}}_{S - \frac{1}{2}}^{0;(1)} - \widehat{\mathbf{G}}_{S + \frac{1}{2}}^{0;(1)}}{\widehat{\mathbf{G}}_{S - \frac{1}{2}}^{0;(3)} - \widehat{\mathbf{G}}_{S + \frac{1}{2}}^{0;(3)}} \neq \frac{0}{0} \; {\rm and} \; \neq \phi \left( \boldsymbol{\mathcal{U}}_S^{0} \right). 
	\end{equation}
	Assuming both $\boldsymbol{\mathcal{U}}_S^{1}(\Delta \tau_0)$ and $\boldsymbol{\mathcal{U}}_S^{0}$ are contained in $\Omega$, then $v \left(\boldsymbol{\mathcal{U}}_S^{1}(\Delta \tau_0) \right) \equiv v_{\min} \equiv v_{\max}$ . Noting that a stable BP numerical scheme should ensure $\boldsymbol{\mathcal{U}}_S^{1}(\Delta \tau) \in \Omega$ for any $\Delta \tau \in (0, \Delta \tau_0)$. Since
	\begin{equation*}
		\phi \left( \boldsymbol{\mathcal{U}}_S^{1}(\Delta \tau) \right) = \frac{\boldsymbol{\mathcal{U}}_S^{1;(1)}(\Delta \tau)}{\boldsymbol{\mathcal{U}}_S^{1;(3)}(\Delta \tau)} = \frac{\boldsymbol{\mathcal{U}}_S^{0;(1)} - \frac{\Delta \tau}{\Delta \xi} \left( \widehat{\mathbf{G}}_{S + \frac{1}{2}}^{0;(1)} - \widehat{\mathbf{G}}_{S - \frac{1}{2}}^{0;(1)} \right)}{\boldsymbol{\mathcal{U}}_S^{0;(3)} - \frac{\Delta \tau}{\Delta \xi} \left( \widehat{\mathbf{G}}_{S + \frac{1}{2}}^{0;(3)} - \widehat{\mathbf{G}}_{S - \frac{1}{2}}^{0;(3)} \right)},
	\end{equation*}
	we emphasize that within the constraint $\Delta \tau > 0$, the sole condition under which 
	\begin{equation*}
		\phi \left( \boldsymbol{\mathcal{U}}_S^{1}(\Delta \tau) \right) = \boldsymbol{\mathcal{U}}_S^{0;(1)} / \boldsymbol{\mathcal{U}}_S^{0;(3)} = \phi \left( \boldsymbol{\mathcal{U}}_S^{0} \right),
	\end{equation*}
	presents a contradiction to \eqref{90}. Therefore, $\phi \left( \boldsymbol{\mathcal{U}}_S^{1}(\Delta \tau_0) \right) \neq \phi \left( \boldsymbol{\mathcal{U}}_S^{0} \right)$. Clearly, $\boldsymbol{\mathcal{U}}_S^{1}(\Delta \tau) $can be expressed as 
	\begin{equation}
	\label{76}
		\begin{aligned}
			\boldsymbol{\mathcal{U}}_S^{1}(\Delta \tau) 
			&= \boldsymbol{\mathcal{U}}_S^0 - \frac{\Delta \tau}{\Delta \xi} \left( \widehat{\mathbf{G}}_{j + \frac{1}{2}}^0 - \widehat{\mathbf{G}}_{j - \frac{1}{2}}^0 \right) \\
			&= \frac{\Delta \tau}{\Delta \tau_0} \left( \boldsymbol{\mathcal{U}}_S^0 - \frac{\Delta \tau_0}{\Delta \xi} \left( \widehat{\mathbf{G}}_{j + \frac{1}{2}}^0 - \widehat{\mathbf{G}}_{j - \frac{1}{2}}^0 \right)\right) 
			+ \left( 1 - \frac{\Delta \tau}{\Delta \tau_0}\right)\boldsymbol{\mathcal{U}}_S^0\\
			&=\frac{\Delta \tau}{\Delta \tau_0}\boldsymbol{\mathcal{U}}_S^{1}(\Delta \tau_0) + \left( 1 - \frac{\Delta \tau}{\Delta \tau_0}\right) \boldsymbol{\mathcal{U}}_S^{0}
		\end{aligned} 
	\end{equation}
	From $\phi \left( \boldsymbol{\mathcal{U}}_S^{1}(\Delta \tau_0) \right) \neq \phi \left( \boldsymbol{\mathcal{U}}_S^{0} \right)$, $v \left( \boldsymbol{\mathcal{U}}_S^{1}(\Delta \tau_0) \right) = v \left( \boldsymbol{\mathcal{U}}_S^{0} \right) = v_{\min} = v_{\max}$, and \eqref{76}, we deduce $v \left(\boldsymbol{\mathcal{U}}^1_S (\Delta \tau) \right) > v_{\max}$ as the last part of the proof of Lemma \ref{778}. Therefore, it is impossible to construct a stable BP numerical scheme when $c \neq v$ for the system \eqref{newsystem}.
\end{proof}

\end{document}
